\clearpage\section{MOS transistor schematics}
Source: \texttt{01\_delay\_elements.sp} and \texttt{02\_taper\_buffer.sp}. Each stage is the MOS inverter below. The first deck compares four independent five-stage chains, with $(W_n,W_p,L)$ in $\mu$m equal to $(1.2,3.6,0.25)$, $(1.2,1.2,0.25)$, $(1.2,3.6,0.5)$ and $(1.2,1.8,0.5)$. Those chains have no added output capacitor.

The drawings use the device and node names of the cited netlists. Equal net labels denote an electrical connection. Unless stated otherwise, NMOS bulk terminals connect to ground and PMOS bulk terminals connect to $V_{DD}$. The bulk wires are omitted from the three-terminal symbols for readability.
\begin{figure}[!htbp]\centering
\resizebox{5.4cm}{!}{\begin{circuitikz}[american,font=\small]

\begin{scope}[shift={(0,0)}]
\node[pmos] (unitp) at (0,1) {};
\node[nmos] (unitn) at (0,-1) {};
\draw (unitp.S) -- (0,2) node[above] {$V_{DD}$};
\draw (unitn.S) -- (0,-2) node[ground] {};
\draw (unitp.D) -- (0,0) -- (unitn.D);
\draw (unitp.G) -- (-1.15,1) -- (-1.15,-1) -- (unitn.G);
\draw (-1.15,0) node[circ] {} -- (-1.8,0) node[above] {a};
\draw (0,0) node[circ] {} -- (0.85,0) node[above] {y};
\node[right] at (0.25,1) {MP};
\node[right] at (0.25,-1) {MN};
\node at (0,2.6) {};
\coordinate (unitin) at (-1.8,0);
\coordinate (unitout) at (0.85,0);
\end{scope}

\end{circuitikz}}
\caption{Two-MOS INV cell used in every delay stage. The named sizing parameters are wn, wp and ll.}
\end{figure}
\begin{figure}[!htbp]\centering
\resizebox{15cm}{!}{\begin{circuitikz}[american,font=\small]
\node[draw,minimum width=1.45cm,minimum height=.8cm] (b0) at (1.3,0) {X1};
\node[draw,minimum width=1.45cm,minimum height=.8cm] (b1) at (3.8,0) {$X2$};
\draw[->] (b0.east) -- (b1.west);
\node[draw,minimum width=1.45cm,minimum height=.8cm] (b2) at (6.3,0) {$X3$};
\draw[->] (b1.east) -- (b2.west);
\node[draw,minimum width=1.45cm,minimum height=.8cm] (b3) at (8.8,0) {$X4$};
\draw[->] (b2.east) -- (b3.west);
\node[draw,minimum width=1.45cm,minimum height=.8cm] (b4) at (11.3,0) {$X5$};
\draw[->] (b3.east) -- (b4.west);
\draw[->] (-.2,0) node[left] {in} -- (b0.west);
\draw[->] (b4.east) -- ++(1,0) node[right] {n5};
\draw (12.7,0) to[C,l=$2\,\mathrm{pF}$] (12.7,-2) node[ground] {};
\end{circuitikz}}
\caption{Tapered five-stage buffer in the second deck: stage $i$ uses $W_n=0.6F^{i-1}$ and $W_p=1.8F^{i-1}\,\mu$m, $L=0.25\,\mu$m; $F=2,3,4,5$. Only n5 drives the 2 pF load.}
\end{figure}
